\documentclass[a4paper,10pt]{article}

\usepackage[utf8]{inputenc}
\usepackage[T1]{fontenc}
\usepackage[top=0.4in, bottom=0.4in, left=0.5in, right=0.5in]{geometry}
\usepackage{enumitem}
\usepackage{titlesec}
\usepackage{hyperref}

\pagestyle{empty}

\titleformat{\section}{\large\bfseries}{}{0em}{}[\titlerule]
\titlespacing*{\section}{0pt}{4pt}{3pt}

\hypersetup{colorlinks=false, pdfborder={0 0 0}}

\begin{document}

%-------------------- HEADER --------------------
\begin{center}
    {\LARGE \textbf{Rupjyoti Barman}} \\[2pt]
    \href{mailto:rupjyoti921@gmail.com}{rupjyoti921@gmail.com} \textbar{} +91-9101171366 \textbar{} \href{https://www.linkedin.com/in/rupjyoti-barman-dev}{linkedin.com/in/rupjyoti-barman-dev} \textbar{} \href{https://github.com/rupjyotibarman}{github.com/rupjyotibarman}
\end{center}

%-------------------- SUMMARY --------------------
\section*{Summary}
\textbf{.NET Backend Developer} with \textbf{3 years} of experience building scalable backend applications, system-level diagnostic platforms, and web solutions for \textbf{HP}. Skilled in \textbf{C\#}, \textbf{ASP.NET Core}, \textbf{Web API}, \textbf{Entity Framework Core}, and \textbf{React}, with experience designing decoupled architectures, integrating native C++ components, and developing enterprise-grade solutions. Proven ability to deliver high-impact features, optimize system performance, and collaborate with global engineering teams across India, Taiwan, and the US.

%-------------------- SKILLS --------------------
\section*{Skills}
\begin{itemize}[leftmargin=*, itemsep=3pt, parsep=0pt, topsep=2pt]
    \item \textbf{Languages:} C\#, C++, SQL, JavaScript, Python
    \item \textbf{Frameworks \& Libraries:} ASP.NET Core (Web API, MVC), Entity Framework Core, WPF, React, Bootstrap, jQuery, LINQ
    \item \textbf{Windows APIs \& Interop:} WMI, WQL, ManagementEventWatcher, WUA API, HPCMSL, DeviceIoControl, P/Invoke, Win32 APIs
    \item \textbf{Architecture \& Design:} Clean Architecture, Layered Architecture, SOLID Principles, OOP, Factory Pattern, Dependency Injection, Controller-Service-Repository Pattern
    \item \textbf{Database, DevOps \& Tools:} SQL Server, PostgreSQL, EF Core (Code-First, Migrations), CI/CD, Azure Pipelines, Git, NuGet, Visual Studio, Postman, Jira, Confluence, Serilog
    \item \textbf{Testing \& Security:} MSTest, Moq, JWT Authentication, RBAC
\end{itemize}

%-------------------- EXPERIENCE --------------------
\section*{Experience}

\textbf{Software Developer, BoldTek India Pvt. Ltd. - Guwahati} \hfill June 2023 - Present \\
\textbf{Client: Hewlett-Packard (HP)}

\begin{itemize}[leftmargin=*, itemsep=5pt, parsep=0pt, topsep=2pt]

\item Engineered 2 production features in an \textbf{ASP.NET Core MVC} application (Controller-Service-Repository, \textbf{EF Core Code-First}, SQL Server, JavaScript, Bootstrap); built an \textbf{ODM self-service dashboard} with 3-tier RBAC (Viewer, ODM, Admin) for INI configuration management, verification submission, and approval tracking with INI key-value validation across \textbf{1000+ HP NPI config files}; developed a \textbf{Release Notes Dashboard} replacing manual Excel-based tracking for the HP Factory Diagnostic team.

\item Built a \textbf{React}-based \textbf{Diagnostic Release Tracking Portal} (POC) using \textbf{React Hooks}, \textbf{React Router}, and reusable components with an \textbf{ASP.NET Core Web API} backend; centralized tracking of diagnostic releases, NuGet package versions, issue resolutions, and development milestones.

\item Led the India team in migrating from a \textbf{shared RPC-based Windows Service architecture} to a \textbf{decoupled three-layer NuGet-based architecture} (Business Logic, C\# Interop Wrapper via P/Invoke, Native C++ DLL) with \textbf{Factory Pattern}-based dependency injection; achieved \textbf{100\% elimination of third-party release delays} (previously 10--15 days) and \textbf{reduced, diagnostic development cycle} from 4 to 3 sprints (40 to 30 working days, \textbf{25\% reduction}) by removing 2 redundant architectural layers, resolving cross-team unhandled exceptions and sleep and hibernate-related service failures in \textbf{HP PC Hardware Diagnostics}.

\item Migrated \textbf{6 hardware diagnostic modules} ( \textbf{Storage:} DRV Read/Verify, SMART Health Monitor, SMART Data Collection; \textbf{Network:} NIC IRQ Verification, ROM/MAC Validation; \textbf{Audio:} Audio Playback Test) onto the new Decoupled Architecture; additionally \textbf{engineered 2 diagnostic modules} from scratch - \textbf{Driver Update Discovery} (WUA API) and \textbf{BIOS Update Check} (HPCMSL, commercial systems) - consolidating all into a \textbf{centralized diagnostic backend} supporting multiple HP client applications with zero external service dependency.

\item Replaced a service-based device discovery architecture (requiring a 2 to 3 minute service restart per hardware change) by \textbf{engineering a real-time WMI monitoring module} (ManagementEventWatcher and WQL on Win32\_PnPEntity) into the shared \textbf{"Common"} library; reduced detection latency \textbf{from 2 to 3 minutes to under 4 seconds} (\textbf{97\% reduction}), adopted across multiple hardware diagnostic APIs.

\item Configured \textbf{Azure Pipelines} to trigger \textbf{NuGet package builds}, publishing the Business Logic package (with dependencies) into the .NET Wrapper layer to produce the final consolidated NuGet package (bundling all layers including Native C++ DLL) for testing and client consumption.

\item Verified diagnostic modules using a \textbf{WPF-based API testing application} for end-to-end UI-level verification alongside \textbf{MSTest} unit test projects across \textbf{Business Logic} and \textbf{.NET Wrapper} layers, leveraging \textbf{Moq} for dependency mocking to ensure reliability prior to client integration.

\item Partnered with cross-functional global teams \textbf{(India, Taiwan, US)} to deliver features, debug production issues, and align on diagnostic API design across release cycles.

\end{itemize}

%-------------------- EDUCATION --------------------
\section*{Education}

\textbf{B.Tech in Computer Science and Engineering} - CGPA: 8.01/10 \hfill May 2023 \\
North Eastern Regional Institute of Science and Technology (NERIST), Arunachal Pradesh

\end{document}